\section*{A Finite Toy Model of the QM--GR Missing Layer}

\paragraph{Purpose.}
This document consolidates the verified framework-driven construction of a
finite toy model for the E018 QM--GR missing-layer residual.  It is a
deliverable for external review, not a root-level completion claim.  Every
claim below is sourced to a per-step artifact from the construction cascade.

\subsection*{1. Setup}

E018 asks for a lawful layer above the QM-side and GR-side access packages.
The re-spined cascade after Step 10 stopped importing the literature's sector
decomposition and instead ran the Six Birds framework on the endpoint access
quotients.  The original obstruction has two relevant faces:

\begin{itemize}
\item P3: route mismatch / scale-descent obstruction, represented natively by
  audit route-consistency and F35 scale descent.
\item P2: problem of time, represented by replacing an external-time dynamic
  with a constrained clock-field relational form.
\end{itemize}

The finite structural carrier used in the native classification is the mode
carrier
\[
L=(d0,d1,d2,d3),
\quad
q_{QM}=(d0,d1,d2),
\quad
q_{GR}=(d0,d2,d3).
\]
The physical-content carrier used later is a finite periodic complex field
\(\psi\) on \(N=8\) sites, plus finite clock registers for relational tests.

\subsection*{2. Native Structural Classification}

\paragraph{Promotion-bridge status.}
Step 17 classified the package \(L\) and the direct-sum package natively in
Foundations-III vocabulary.  \(B_{QM\to L}\) passes the eight gates and has
status \(\texttt{candidate}\) on the toy.  The direct-sum package is also
strict over QM, so strictness is not the discriminator.  It fails
\(\texttt{G\_nosmuggle}\) because it duplicates the shared observables \(d0\)
and \(d2\).  Source:
\texttt{steps/step17\_native\_promotion\_bridge\_classification\_artifacts/step17\_results\_summary.md}.

\paragraph{F37 joint quotient.}
Step 18 applied F37.  With
\[
J=O_L,\quad j=\pi_L,\quad a:J\to O_{QM},\quad b:J\to O_{GR},
\]
the commuting residuals are \(a\circ j=q_{QM}\) and \(b\circ j=q_{GR}\), both
with residual \(0.0\).  Since \(L\) is admissible while the product/direct-sum
control is not, \(q_{QM}\) and \(q_{GR}\) are not complementary on the toy.
Source:
\texttt{steps/step18\_f37\_complementarity\_f24\_resolution\_artifacts/step18\_results\_summary.md}.

\paragraph{F24 resolution family.}
Step 20 built operational selectors for the F24 resolution families, and
Step 21 rebuilt them structurally from actual maps rather than selector tags.
The candidate package \(L\) recomputes uniquely as
\(\texttt{BridgeMediatedRole}\).  The discrimination suite covers the eight
families exactly once.  Sources:
\texttt{steps/step20\_f24\_predicate\_construction\_artifacts/step20\_results\_summary.md}
and
\texttt{steps/step21\_f24\_structural\_predicates\_artifacts/step21\_results\_summary.md}.

\paragraph{F35 scale descent.}
Step 19 applied F35 to the audited joint law on \(L\).  The joint law descends
and renormalizes across \(Q1\to Q2\) and \(Q2\to Q3\).  The non-descending
control fails first by coarsening-fiber violation and then by a nonempty
obstruction set.  Source:
\texttt{steps/step19\_f35\_scale\_descent\_renormalization\_artifacts/step19\_results\_summary.md}.

\paragraph{F51 common-refinement unification.}
Step 22 applied F51.  \(L\) is a common-refinement unification of \(q_{QM}\)
and \(q_{GR}\) on the finite carrier: the parent-to-child projection residuals
are \(0.0\), \(L\) strictly refines both children, and child descent statuses
are compatible.  The status-conflict control fails.  Source:
\texttt{steps/step22\_f51\_unification\_common\_refinement\_artifacts/step22\_results\_summary.md}.

\paragraph{Durable structural constraints.}
The cascade also produced reusable constraints:

\begin{itemize}
\item A lawful package is not the direct-sum package; the native distinction is
  \(\texttt{G\_nosmuggle}\), not strictness alone.  Source: Step 17.
\item A shared currency must compress rather than staple endpoint currencies:
  \(d_u=4<d_{union}=6\).  Source: Step 13.
\item Refinement must be a morphism of the two-access diagram, commuting with
  both endpoint readouts.  Source: Step 14.
\item A single audit exists iff endpoint audits agree on shared modes; this
  recovers P3 as route-consistency on the overlap.  Source: Step 15.
\end{itemize}

These are native/toy structural constraints.  They are not claims about all
possible physical models.

\subsection*{3. Physical Toy Content}

\paragraph{Derived audits and sourcing.}
Step 23 first tested compatible quadratic audit forms.  Step 24 then replaced
chosen overlap agreement by derived field audits.  From the same finite field
\(\psi\),
\[
\rho=|\psi|^2,\quad
j=\operatorname{Im}(\overline{\psi}\nabla\psi),
\]
while
\[
T_{00}=|\nabla\psi|^2+V|\psi|^2,\quad
T_{0i}=\operatorname{Re}(\overline{\nabla\psi}\,\Delta\psi).
\]
The Born and stress-energy audits do not cohere as one shared audit: the
density proportional residual is \(0.5207123913178113\), and the transport
proportional residual is \(0.9994686548331937\).  The positive relation is
sourcing: \(T[\psi]\) is a deterministic functional of the shared field, with
direct sourcing residual \(0.0\).  Sources: Steps 23 and 24.

\paragraph{Co-sourcing package.}
Step 25 built \(L\) as a field carrier.  The QM descent reads
\(\text{Born}[\psi]\), while the GR descent reads \(T[\psi]\).  The co-sourced
package has max Born residual \(0.0\) and max stress-energy residual \(0.0\);
the non-co-sourced control fails with max stress-energy residual
\(0.3806026603775931\).  Source:
\texttt{steps/step25\_sourcing\_unification\_artifacts/step25\_results\_summary.md}.

\paragraph{Semiclassical dynamics.}
Step 26 added a finite self-sourced field-potential dynamic:
\[
V[\psi]=B+\kappa T_{00}[\psi],
\qquad
H[V]=\frac12\Delta_{per}+\operatorname{diag}(V).
\]
For moderate coupling \(\kappa=0.3\), the Hartree-like iteration converges to
a self-consistent stationary field-potential point with fixed-point residual
\(3.4936318023019015\cdot 10^{-16}\) and eigen residual
\(6.377748916558787\cdot 10^{-16}\).  The over-strong control fails with
fixed-point residual \(1.4285712942487314\).  Source:
\texttt{steps/step26\_semiclassical\_dynamics\_artifacts/step26\_results\_summary.md}.

\paragraph{Timeless and relational form.}
Step 27 represented the stationary Step 26 dynamic as a constrained
clock-field state.  The compatible constraint has kernel dimension \(1\) at
tolerance \(10^{-8}\), \(\|C\Psi\|=8.547035876629756\cdot 10^{-16}\), and
clock conditioning recovers the stationary dynamic with max residual
\(9.79195701998934\cdot 10^{-16}\).  The empty-kernel control fails.  Source:
\texttt{steps/step27\_constraint\_problem\_of\_time\_artifacts/step27\_results\_summary.md}.

Step 28 extended this to one non-stationary history.  A non-eigenstate field
evolves with max phase-invariant state change \(0.18847313451682898\), and
the sourced potential changes with max norm \(0.001673580612625418\).  The
generated relational constraint has \(\|C\Psi\|=4.220196766037856\cdot
10^{-16}\), and the wrong-generator control fails with residual
\(0.0925307478333232\).  Source:
\texttt{steps/step28\_nonstationary\_relational\_histories\_artifacts/step28\_results\_summary.md}.

\paragraph{Recovered known structure.}
At toy fidelity, the physical content recovers two familiar structures:
\begin{itemize}
\item semiclassical sourcing, where a field determines stress-energy content;
\item relational clock readout, where global clock-field states are constrained
  and dynamics appears through conditioning.
\end{itemize}
This is a recovery of known structural patterns in a finite carrier, not an
empirical novelty claim.

\subsection*{4. Grade and Scope}

\paragraph{Grade.}
The native classifications are graded \(\texttt{native-law-verified}\) in the
limited sense that the construction applies the framework laws and validators
from the repository to finite carriers.  The physical content is graded
\(\texttt{toy-diagnostic}\): finite fields, finite clocks, chosen couplings,
and explicit controls.

\paragraph{Can-fail controls.}
The construction uses controls that fail: direct-sum package, route-mismatch
audit, non-descending law, non-co-sourced field, over-strong dynamics, empty
kernel, and wrong relational generator.

\paragraph{Posited or adapted parts.}
The following components are honest toy choices:
\begin{itemize}
\item Step 23 uses chosen quadratic audit forms before Step 24 derives audits
  from a field.
\item Step 26 uses a finite Hamiltonian and toy stress-energy density.
\item Step 27 uses a compatible finite clock for the stationary history.
\item Step 28 uses a history-adapted open-chain constraint after generating
  the self-sourced sequence.
\end{itemize}

\paragraph{Scope boundary.}
The consolidated result is a native-verified structural classification plus a
toy-diagnostic recovery of semiclassical sourcing and relational time.  It is
not a complete physical model, not an empirical prediction, and not a
frame-transfer certificate.  Frame transfer beyond the finite carriers is an
external-review task.

\subsection*{5. Independently Checkable Content}

An external reviewer can verify:
\begin{itemize}
\item native status and gate records by rerunning Step 17;
\item F37/F24/F35/F51 computations by rerunning Steps 18, 19, 21, and 22;
\item derived sourcing by rerunning Steps 24 and 25;
\item finite self-consistency by rerunning Step 26;
\item stationary and non-stationary relational constraints by rerunning Steps
  27 and 28;
\item Step 29's claim ledger in
  \texttt{content\_classification\_step29.csv}.
\end{itemize}

\subsection*{6. Open Obligations}

\begin{itemize}
\item Replace the Step 28 history-adapted constraint with a less adapted
  nonlinear constraint tested over broader history families.
\item Enrich the field carrier toward more faithful QM content: Hilbert-space
  observables, measurement records, and stability of record readouts.
\item Enrich the geometric side toward more faithful GR content: Lorentzian
  signature, diffeomorphism-style redundancy, and constraint algebra tests.
\item Test robustness under different lattices, couplings, dimensions, and
  clock constructions.
\item Perform the real physics frame-transfer review.
\end{itemize}

\paragraph{Consolidated verdict.}
Within the declared finite carriers and verified artifacts, \(L\) is a
candidate common-refinement package that is native-classified as admissible,
BridgeMediatedRole, scale-consistent, and F51-unifying on the toy.  With
finite physical content, \(L\) carries a field whose Born and stress-energy
descents are co-sourced; it admits self-consistent semiclassical dynamics and
stationary plus non-stationary relational clock-field representations.  This
is the construction product for external review.
